\documentclass[11pt,a4paper]{article}
\usepackage[T1]{fontenc}
\usepackage{lmodern}
\usepackage[margin=27mm]{geometry}
\usepackage{amsmath,amssymb,amsthm,mathtools}
\usepackage[colorlinks=true,linkcolor=blue,urlcolor=blue,citecolor=blue]{hyperref}
\usepackage{microtype}
\newtheorem{theorem}{Theorem}
\newtheorem{lemma}{Lemma}
\newcommand{\R}{\mathbb R}
\newcommand{\D}{\mathcal D}
\newcommand{\Lop}{\mathcal L}
\newcommand{\Kop}{\mathcal K}
\newcommand{\Tcal}{\mathcal T}
\newcommand{\Ical}{\mathcal I}
\newcommand{\norm}[1]{\left\lVert#1\right\rVert}
\newcommand{\ip}[2]{\left\langle#1,#2\right\rangle}
\newcommand{\mathscr}{\mathcal}
\title{An embedded eigenvalue in an asymptotically flat\protect\\purely twisted strip}
\author{Siavash Sadeghi}
\date{2 November 2026; revised 3 November 2026}
\hypersetup{pdftitle={An embedded eigenvalue in an asymptotically flat purely twisted strip},pdfauthor={Siavash Sadeghi},pdfsubject={Proposed counterexample for AIM 1007}}
\begin{document}
\maketitle
\begin{abstract}
We give a counterexample to the universal absence assertion in AIM problem 1007. An explicit one-parameter family of positive twisting profiles contains an embedded eigenvalue for every sufficiently large value of a scale parameter. The proof treats the complete transverse operator, uses a contraction in global Sobolev spaces, and enforces the longitudinal matching condition through a continuous antiperiodic phase family.
\end{abstract}

\noindent\textbf{Target and status.}
This is a proposed negative resolution of
\href{https://github.com/MColbrook/AIM/blob/ee5c5c52f527db2f310cb0734bbebc9f7b95aac8/problems/1007-embedded-eigenvalues-in-asymptotically-flat-twisted-strips.md}{AIM problem 1007},
at repository revision \texttt{ee5c5c52f527}. It concerns the
point-spectrum question strictly above the first transverse threshold;
the separate resonance question in the source is outside this target.
The accompanying audit records a fresh AI review and symbolic checks,
not a human peer review or a formal verification.

For $a>0$ let $\Omega_a=\R\times(-a,a)$ and
$E_j=(j\pi/(2a))^2$. The admissible class consists of
$q\in L^\infty(\R)$ such that $q\ge0$, $q$ is not almost everywhere zero,
$a\norm q_\infty\le\sqrt2$, and
$\operatorname*{ess\,sup}_{|s|>R}q(s)\to0$ as $R\to\infty$.
With $f_q(s,t)=\sqrt{1+t^2q(s)^2}$, let $H_q$ be the selfadjoint operator
in $L^2(\Omega_a,f_q\,ds\,dt)$ defined by
\[
 h_q[\psi]=\int_{\Omega_a}
   \bigl(f_q^{-1}|\psi_s|^2+f_q|\psi_t|^2\bigr)\,ds\,dt,
 \qquad D(h_q)=H_0^1(\Omega_a).
\]
We denote this admissible class by $\mathscr Q_a$.

\begin{theorem}\label{thm:main}
Set
\[
 m=\frac13-\frac2{\pi^2},\qquad
 k=\frac1{\sqrt{2m}},\qquad
 \lambda_*=\frac{\pi^2}{4}+k^2=\frac{\pi^4}{4(\pi^2-6)}.
\]
There is a finite $S_0$ such that, for every $S\ge S_0$, some
$\delta\in[0,2\pi]$ has the following property. The profile
\begin{equation}\label{eq:profile}
 q_{S,\delta}(s)
 =\left[\frac{16k}{S+|s|}
     \left(1+\frac12\cos(2k|s|+\delta)\right)\right]^{1/2}
\end{equation}
belongs to $\mathscr Q_1$, and $H_{q_{S,\delta}}$ has a nonzero eigenfunction
in $H^2(\R\times(-1,1))\cap H_0^1(\R\times(-1,1))$ at energy $\lambda_*$.
In particular, $E_1<\lambda_*<E_2$.
For every $a>0$, the profile $q_a(s)=a^{-1}q_{S,\delta}(s/a)$ belongs to
$\mathscr Q_a$ and has an embedded eigenvalue $a^{-2}\lambda_*$.
\end{theorem}

Only the phase is selected by existence. It is specified below as a zero
of a continuous real function obtained from convergent linear inverses.
No transverse truncation or numerical spectral assertion is used.

The intrinsic Dirichlet Laplace--Beltrami operator describes a particle
on a two-dimensional twisted ribbon, confined transversely by its edges
but free to propagate longitudinally. This is a surface or ribbon
quantum waveguide; the result concerns this model rather than a solid
three-dimensional twisted tube. In the asymptotic free guide, the
transverse mode $\chi_j$ has dispersion $\lambda=E_j+p^2$, where $p$
is longitudinal momentum. A channel is open when $\lambda>E_j$,
permitting longitudinal propagation. Since $\lambda_*\in(E_1,E_2)$,
the first channel is open, and the constructed state's component
$u\chi_1$ in that channel is nonzero. Nevertheless, its eigenfunction
is square integrable in $f_q\,ds\,dt$ and defines a normalizable
stationary state: a bound state in the continuum. Localization concerns
the probability distribution, not compact support. Under Schr\"odinger
evolution, the eigenstate changes only by a global phase, so its
probability density remains stationary despite the available
longitudinal escape channel. The construction achieves this through
oscillatory long-range twisting.

\section{Half-strip formulation and uniform estimates}

Work with $a=1$. For $J=(-1,1)$ define the transverse Hilbert space
and Dirichlet operator by
\[
 \mathfrak h=L^2(J,dt),\qquad T=-\partial_t^2,\qquad
 D(T)=H^2(J)\cap H_0^1(J).
\]
Its normalized orthonormal eigenbasis $(\chi_j)_{j\ge1}$ is
\[
 \chi_j(t)=
 \begin{cases}
  \cos(j\pi t/2),&j\text{ odd},\\
  \sin(j\pi t/2),&j\text{ even},
 \end{cases}
 \qquad T\chi_j=E_j\chi_j,\qquad E_j=(j\pi/2)^2.
\]
Write $\chi=\chi_1=\cos(\pi t/2)$. The transverse orthogonal projections are
\[
 P_Tw=\chi_1\int_J\overline{\chi_1(t)}w(t)\,dt,\qquad
 Q_T=I_{\mathfrak h}-P_T.
\]
Their closed-channel space and the transverse restriction are
\[
 \mathfrak h_Q=\operatorname{Ran}Q_T
 =\{\chi_1\}^{\perp}
 =\overline{\operatorname{span}}\{\chi_j:j\ge2\},\qquad
 T_Q=T\big|_{D(T)\cap\mathfrak h_Q}.
\]
The space $\mathfrak h_Q$ reduces $T$, so $T_Q$ is selfadjoint in
$\mathfrak h_Q$ and $T_Q\ge E_2$. All modes $j\ge2$ are retained.
Since $\pi^2>8$, the formula for $\lambda_*$ gives $\lambda_*<E_2$.

On the half-strip and its unweighted auxiliary Hilbert space,
\[
 \D=(0,\infty)\times J,\qquad
 \mathfrak H=L^2(\D,ds\,dt)=L^2((0,\infty);\mathfrak h),
\]
define the fibrewise extensions and their closed-channel range by
\[
 P_{\D}=I_{L^2(0,\infty)}\otimes P_T,\qquad
 Q_{\D}=I_{L^2(0,\infty)}\otimes Q_T,\qquad
 \mathfrak H_Q=\operatorname{Ran}Q_{\D}
              =L^2((0,\infty);\mathfrak h_Q).
\]
These projections are orthogonal in the unweighted auxiliary spaces
used for the divergence-form equation, not in the physical space
$L^2(\D,f_\delta\,ds\,dt)$. The perturbed operator need not preserve
this decomposition.

Put $\rho(s)=S+s$, with $S\ge1$. Constants below may depend on the
fixed $k$ and $\lambda_*$, but are independent of $\delta$ and sufficiently
large $S$. For $\mathcal U=(0,\infty)$ or $\D$ and $j\ge0$ define
\[
 Y_j(\mathcal U)=\{F:\rho^jF\in L^2(\mathcal U)\},\qquad
 X_j(\mathcal U)=
 \{v:\rho^j\partial^\nu v\in L^2(\mathcal U),\ |\nu|\le2\},
\]
with the norms
\[
 \norm{F}_{Y_j(\mathcal U)}=\norm{\rho^jF}_{L^2(\mathcal U)},\qquad
 \norm{v}_{X_j(\mathcal U)}
 =\sum_{|\nu|\le2}\norm{\rho^j\partial^\nu v}_{L^2(\mathcal U)}.
\]
All underlying norms use unweighted Lebesgue measure. On the half-line,
$\partial^\nu$ denotes the longitudinal derivatives of orders $0,1,2$;
on the half-strip it denotes all derivatives in $s,t$ of total order
at most two. Thus $X_0(\mathcal U)=H^2(\mathcal U)$.
Projection and boundary conditions are additional conditions, not
part of these definitions. For $\R^2$-valued functions we write
$X_j((0,\infty);\R^2)$ and $Y_j((0,\infty);\R^2)$, using the same
definitions with Euclidean norms. All constructions below preserve
real-valued functions; we use their real subspaces when selecting the phase.

On $s\ge0$ set
\begin{equation}\label{eq:b}
 b_\delta(s)=\frac{16k}{\rho(s)}
       \left(1+\frac12\cos(2ks+\delta)\right),\qquad
 f_\delta(s,t)=\sqrt{1+t^2b_\delta(s)}.
\end{equation}
For each fixed derivative order $j$,
\begin{equation}\label{eq:b-bound}
 |b_\delta^{(j)}(s)|\le C_j\rho(s)^{-1}.
\end{equation}
The eigenvalue equation for the stated form is exactly $\Lop_\delta\psi=0$, where
\begin{align}
 \Lop_\delta
 &=-\partial_s(f_\delta^{-1}\partial_s)
   -\partial_t(f_\delta\partial_t)-\lambda_*f_\delta,
 \label{eq:L}\\
 \Lop_{\rm free}&=-\partial_s^2-\partial_t^2-\lambda_*,
 \qquad \Kop_\delta=\Lop_\delta-\Lop_{\rm free}.\notag
\end{align}
On expanding derivatives, every coefficient of the order-two differential
operator $\Kop_\delta$ is bounded by $C/\rho(s)$. They are combinations of
$f_\delta^{-1}-1$, $\partial_sf_\delta^{-1}$, $f_\delta-1$,
$\partial_tf_\delta$, and $\lambda_*(f_\delta-1)$. Consequently
\begin{align}
 \norm{\Kop_\delta v}_{Y_1(\D)}&\le C\norm{v}_{H^2(\D)},\label{eq:K1}\\
 \norm{\Kop_\delta v}_{Y_1(\D)}&\le \frac CS\norm{v}_{X_1(\D)},
 &\norm{\Kop_\delta v}_{Y_2(\D)}&\le C\norm{v}_{X_1(\D)}.
 \label{eq:K2}
\end{align}
These statements do not require a longitudinal boundary condition on $v$.
The coefficient maps in these operator norms are continuous and
$2\pi$-periodic in $\delta$.

Let $N$ be the longitudinal Neumann operator in $L^2(0,\infty)$:
\[
 N=-\frac{d^2}{ds^2},\qquad
 D(N)=\{u\in H^2(0,\infty):u'(0)=0\}.
\]

\begin{lemma}[Free closed inverse]\label{lem:closed-free}
The operator
\[
 A_0=N\otimes I_{\mathfrak h_Q}
       +I_{L^2(0,\infty)}\otimes T_Q-\lambda_*I
 \quad\text{in }\mathfrak H_Q
\]
has domain
\[
 D(A_0)=\{v\in H^2(\D)\cap\mathfrak H_Q:
          v(s,\pm1)=0,\ \partial_sv(0,t)=0\},
\]
with boundary conditions in the Sobolev trace sense. Its inverse satisfies
\[
 \norm{A_0^{-1}F}_{H^2(\D)}\le C\norm{F}_{L^2(\D)},\qquad
 \norm{A_0^{-1}F}_{X_1(\D)}\le C\norm{F}_{Y_1(\D)},
\]
for $F\in\mathfrak H_Q$, with the additional weighted condition in the
second estimate, and with $C$ independent of $S\ge1$.
\end{lemma}

\begin{proof}
Expand $v(s,t)=\sum_{j\ge2}v_j(s)\chi_j(t)$ in the transverse
Dirichlet eigenbasis. The $j$-th component of $A_0v$ is
$-v_j''+(E_j-\lambda_*)v_j$. Applying the continuous cosine transform
in the longitudinal variable $s$, corresponding to Neumann data at zero,
gives the multipliers
\[
 A_0:\ \xi^2+E_j-\lambda_*,
 \qquad A_0^{-1}:\ (\xi^2+E_j-\lambda_*)^{-1}.
\]
Here $E_j$ comes from the discrete transverse expansion and $\xi^2$
from the continuous longitudinal transform. The spectral gap gives
\[
 \sup_{j\ge2,\,\xi\ge0}
 \frac{1+\xi^2+E_j}{\xi^2+E_j-\lambda_*}<\infty.
\]
This proves the unweighted estimate, including mixed derivatives, and
identifies the domain as the stated closed-channel subspace of $H^2(\D)$.

For the weighted estimate use $w(s)=S+\sqrt{1+s^2}$.
Then $\rho\le w\le2\rho$, $w'(0)=0$, and $|w'|+|w''|\le C$.
Choose bounded smooth truncations $w_N$ tending pointwise to $w$, with
$w_N\le w$, $w_N'(0)=0$, and uniformly bounded first and second
derivatives. Such truncations result by smoothly flattening
$\sqrt{1+s^2}$ above level $N$ before adding $S$.
For $v=A_0^{-1}F$, multiplication by $w_N$ preserves both boundary
conditions and $\mathfrak H_Q$. Hence
\[
 A_0(w_Nv)=w_NF-2w_N'v_s-w_N''v.
\]
The unweighted estimate gives
\[
 \norm{w_Nv}_{H^2(\D)}
 \le C\bigl(\norm{wF}_{L^2(\D)}+\norm{v}_{H^2(\D)}\bigr)
 \le C\norm{wF}_{L^2(\D)}.
\]
The product rule bounds all $w_N\partial^\nu v$ for $|\nu|\le2$
by the same right-hand side. Fatou's lemma proves the assertion.
The weight chosen here preserves the Neumann condition at zero.
\end{proof}

\begin{lemma}[Exact elimination of the closed modes]\label{lem:eliminate}
For sufficiently large $S$ and every scalar $u\in H^2(0,\infty)$, there
is a unique $z=z_\delta[u]\in X_1(\D)\cap\mathfrak H_Q$, with transverse Dirichlet
conditions and $z_s(0,t)=0$ such that
\[
 Q_{\D}\Lop_\delta(u\chi+z)=0\quad\text{on }\D.
\]
It satisfies
\begin{equation}\label{eq:z-est}
 \begin{aligned}
 \norm{z_\delta[u]}_{X_1(\D)}&\le C\norm{u}_{H^2(0,\infty)},\\
 \norm{\ip{\chi}{\Kop_\delta z_\delta[u]}_{\mathfrak h}}_{Y_2((0,\infty))}
 &\le C\norm{u}_{H^2(0,\infty)}.
 \end{aligned}
\end{equation}
The resulting linear maps are continuous and $2\pi$-periodic in $\delta$.
\end{lemma}

\begin{proof}
Let $R_0=A_0^{-1}$. The equation for $z$ is the fixed-point equation
\begin{equation}\label{eq:z-fp}
 z=-R_0Q_{\D}\Kop_\delta(u\chi)-R_0Q_{\D}\Kop_\delta z.
\end{equation}
On the closed subspace of $X_1(\D)\cap\mathfrak H_Q$ carrying the
stated boundary conditions, the second operator has norm at most $C/S$,
by Lemma~\ref{lem:closed-free} and \eqref{eq:K2}. The first term has
$X_1(\D)$ norm at most $C\norm{u}_{H^2(0,\infty)}$ by \eqref{eq:K1}.
The Neumann series proves existence, uniqueness and the first bound.
The second follows from the last estimate in \eqref{eq:K2}.
The lifted projection $Q_{\D}$ is bounded on the strip spaces used:
it acts fibrewise in $t$, and the complementary transverse projection
$P_T$ has the fixed smooth range spanned by $\chi$.
Uniform convergence of the Neumann series proves parameter continuity.

No boundary condition on $u$ is imposed here. This is an interior
differential elimination with $z_s(0)=0$. The full longitudinal boundary
condition will be enforced later by $u'(0)=0$.
\end{proof}

\section{The exact scalar equation}

Define the scalar functions
\begin{align}
 A(b)&=\int_J(1+t^2b)^{-1/2}\chi^2\,dt,\notag\\
 B(b)&=\int_J(1+t^2b)^{1/2}
           (\chi'^2-\lambda_*\chi^2)\,dt,\notag\\
 \alpha(s)&=A(b_\delta(s)),\qquad
 \beta(s)=B(b_\delta(s)).\label{eq:alpha-beta}
\end{align}
Writing $\psi=u\chi+z_\delta[u]$, its remaining equation is precisely
\begin{equation}\label{eq:scalar-u}
 -(\alpha u')'+\beta u+F_\delta u=0,\qquad
 F_\delta u=\ip{\chi}{\Kop_\delta z_\delta[u]}_{\mathfrak h}.
\end{equation}
Here the inner product is taken in $t$, leaving a scalar function of $s$.
Indeed, $P_{\D}\Lop_{\rm free}z=0$, and transverse integration by parts gives
the coefficient $\beta$ in \eqref{eq:scalar-u}.

The elementary identities
\[
 \int_Jt^2\chi^2=m,\qquad
 \int_Jt^2\chi'^2=E_1m+1
\]
give, uniformly for small nonnegative $b$,
\begin{equation}\label{eq:expansion}
 A(b)=1-\frac m2b+O(b^2),\qquad
 B(b)=-k^2+\frac{1-mk^2}{2}b+O(b^2).
\end{equation}
The second identity follows by integrating
$\int t^2\chi(-\chi'')=E_1m$ once and using $\int 2t\chi\chi'=-1$.

Put $y=\sqrt\alpha\,u$. Direct differentiation gives
\[
 -(\alpha u')'+\beta u
 =\sqrt\alpha\,(-y''-W_\delta y),\qquad
 W_\delta=-\frac\beta\alpha-\frac{\alpha''}{2\alpha}
                       +\frac{(\alpha')^2}{4\alpha^2}.
\]
Thus \eqref{eq:scalar-u} is equivalent to
\begin{equation}\label{eq:y-exact}
 y''+W_\delta y
 =\alpha^{-1/2}F_\delta(\alpha^{-1/2}y).
\end{equation}
Taylor's formula for the smooth functions $A,B$, and \eqref{eq:b-bound},
give
\begin{equation}\label{eq:Omega-expand}
 W_\delta
 =k^2+(mk^2-\tfrac12)b_\delta+\frac m4b_\delta''+r_{1,\delta},
 \qquad |r_{1,\delta}(s)|\le C\rho(s)^{-2}.
\end{equation}
All remainders here are bounded by a constant times
$b^2+|bb''|+(b')^2$; no uncontrolled longitudinal derivative occurs.
The choice of $k$ gives $mk^2=1/2$. Differentiating \eqref{eq:b} twice
therefore yields
\begin{equation}\label{eq:resonance}
 W_\delta=k^2-\frac{4k}{\rho}\cos(2ks+\delta)+r_\delta,
 \qquad |r_\delta(s)|\le C\rho(s)^{-2}.
\end{equation}
The complete equation, with every closed mode included, is
\begin{equation}\label{eq:full-scalar}
 \mathfrak L_\delta y:=
 y''+\left(k^2-\frac{4k}{\rho}\cos(2ks+\delta)\right)y
 =\Tcal_\delta y,
\end{equation}
where
\[
 \Tcal_\delta y=
 \alpha^{-1/2}F_\delta(\alpha^{-1/2}y)-r_\delta y.
\]
Since $\alpha$ stays bounded away from zero, its inverse square root
and its first two derivatives are uniformly bounded.
Lemma~\ref{lem:eliminate} and \eqref{eq:resonance} show that
\begin{equation}\label{eq:T-bound}
 \Tcal_\delta:H^2(0,\infty)\longrightarrow Y_2((0,\infty)),\qquad
 \norm{\Tcal_\delta y}_{Y_2((0,\infty))}\le C\norm{y}_{H^2(0,\infty)}.
\end{equation}
This is an exact bounded linear operator, generally nonlocal in $s$.
It is real, continuous and $2\pi$-periodic in $\delta$ in the indicated
operator norm. Uniform bounds for the Taylor remainders and their
$\delta$ derivatives give the same assertion for multiplication by
$r_\delta$.

\section{A scalar right inverse with a small norm}

\begin{lemma}\label{lem:scalar}
For sufficiently large $S$ there are real functions
$h_\delta\in H^2(0,\infty)$ and real bounded linear maps
$G_\delta:Y_2((0,\infty))\to H^2(0,\infty)$ such that
\[
 \mathfrak L_\delta h_\delta=0,\qquad
 \mathfrak L_\delta G_\delta g=g,\qquad
 \norm{G_\delta}_{Y_2((0,\infty))\to H^2(0,\infty)}\le C/S.
\]
The family $G_\delta$ is continuous and $2\pi$-periodic, while
$h_\delta$ is continuous and satisfies $h_{\delta+2\pi}=-h_\delta$.
For the vector $e_\delta\in\R^2$ and the scalar trace functional
$\ell_\delta:H^2(0,\infty)\to\R$ defined by
\[
 e_\delta=
 \begin{pmatrix}\cos(3\pi/4-\delta/2)\\
                 \sin(3\pi/4-\delta/2)\end{pmatrix},
 \qquad
 \ell_\delta(y)=e_\delta^{\mathsf T}
            \begin{pmatrix}y(0)\\y'(0)/k\end{pmatrix},
\]
the normalization is
\begin{equation}\label{eq:normalization}
 \ell_\delta(h_\delta)=1/S,\qquad
 \ell_\delta(G_\delta g)=0.
\end{equation}
\end{lemma}

\begin{proof}
Start with $\mathfrak L_\delta y=g$, with scalar
$g\in Y_2((0,\infty))$, and put $v=(y,y'/k)^{\mathsf T}\in\R^2$.
Let $R(s),E$ be the following real $2\times2$ matrices and $e_2\in\R^2$:
\[
 R(s)=\begin{pmatrix}\cos ks&\sin ks\\-\sin ks&\cos ks\end{pmatrix},
 \quad E=\begin{pmatrix}0&0\\1&0\end{pmatrix},
 \quad e_2=\begin{pmatrix}0\\1\end{pmatrix}.
\]
Writing $v=R(s)a$, with $a(s)\in\R^2$, converts the equation exactly to
\begin{equation}\label{eq:rotate}
 a'=\frac{M_\delta+N_\delta(s)}\rho a+R(s)^{-1}e_2g/k,\qquad
 M_\delta=\begin{pmatrix}\sin\delta&\cos\delta\\
                         \cos\delta&-\sin\delta\end{pmatrix}.
\end{equation}
To verify this, $M_\delta+N_\delta(s)$ is the matrix
$4\cos(2ks+\delta)R(s)^{-1}ER(s)$.
Its mean over the period $\pi/k$ is $M_\delta$, so $N_\delta$ is
bounded, periodic and has mean zero. Also $M_\delta^2=I$ and
$M_\delta e_\delta=-e_\delta$.

Define $Z_\delta(s)=\int_0^sN_\delta(r)\,dr$.
It is bounded and periodic, with $Z_\delta(0)=0$.
For sufficiently large $S$, $H_\delta(s)=I+Z_\delta(s)/\rho(s)$ is
uniformly invertible. These are real $2\times2$ matrices.
Set $a=H_\delta x$, with $x(s)\in\R^2$. The exact transformed equation is
\begin{equation}\label{eq:normal}
 x'=\frac{M_\delta}\rho x+D_\delta(s)x+B_\delta(s)g,
\end{equation}
where
\begin{align*}
 D_\delta(s)&=\rho^{-2}H_\delta^{-1}
 \bigl([M_\delta,Z_\delta]+N_\delta Z_\delta+Z_\delta\bigr),\\
 B_\delta(s)&=H_\delta^{-1}R(s)^{-1}e_2/k.
\end{align*}
In particular $|D_\delta(s)|\le C\rho^{-2}$ and $|B_\delta(s)|\le C$.
Here $D_\delta(s)\in\R^{2\times2}$ and $B_\delta(s)\in\R^2$;
both are continuous and $2\pi$-periodic in $\delta$.

The orthogonal projections $\Pi_\pm=(I_{\R^2}\pm M_\delta)/2$ act on
$\R^2$, independently of the transverse and half-strip projections.
For $F\in Y_2((0,\infty);\R^2)$ define
\begin{equation}\label{eq:I}
 (\Ical_\delta F)(s)
 =\rho(s)^{-1}\int_0^s\rho(r)\Pi_-F(r)\,dr
  -\rho(s)\int_s^\infty\rho(r)^{-1}\Pi_+F(r)\,dr.
\end{equation}
This solves $x'=M_\delta x/\rho+F$ and has $\Pi_-x(0)=0$. Moreover,
\begin{equation}\label{eq:I-bound}
 \norm{\Ical_\delta F}_{L^2((0,\infty);\R^2)}
 \le\frac CS\norm{F}_{Y_2((0,\infty);\R^2)}.
\end{equation}
Indeed, writing $f=\rho^2F\in L^2((0,\infty);\R^2)$, the stable kernel is
$1_{r\le s}/(\rho(s)\rho(r))$, whose squared Hilbert--Schmidt norm is
\[
 \int_0^\infty\rho(r)^{-2}\int_r^\infty\rho(s)^{-2}\,ds\,dr
 =\frac1{2S^2}.
\]
The unstable kernel is $1_{r\ge s}\rho(s)/\rho(r)^3$.
Its row integrals are $1/(2\rho(s))$; its column integrals are
$(\rho(r)^2-S^2)/(2\rho(r)^3)$. Both are at most $1/(2S)$.
The Schur estimate proves \eqref{eq:I-bound}.

Multiplication by $D_\delta$ maps $L^2((0,\infty);\R^2)$ to
$Y_2((0,\infty);\R^2)$ with uniformly bounded norm.
Hence $I-\Ical_\delta D_\delta$ has a convergent Neumann inverse on
$L^2((0,\infty);\R^2)$ for large $S$. Define
\begin{align*}
 x_{h,\delta}&=(I-\Ical_\delta D_\delta)^{-1}(\rho^{-1}e_\delta),\\
 x_{g,\delta}&=(I-\Ical_\delta D_\delta)^{-1}
                         \Ical_\delta(B_\delta g).
\end{align*}
These solve \eqref{eq:normal} with $g=0$ and with the given $g$,
respectively. They are in $H^1((0,\infty);\R^2)$:
differentiating \eqref{eq:I} in the distributional sense gives
\eqref{eq:normal}, whose right-hand side is in $L^2((0,\infty);\R^2)$. The integral formulas also give the trace data
\[
 \Pi_-x_{h,\delta}(0)=e_\delta/S,\qquad
 \Pi_-x_{g,\delta}(0)=0.
\]
Define $h_\delta$ and $G_\delta g$ as the first components of
$RH_\delta x_{h,\delta}$ and $RH_\delta x_{g,\delta}$.
Since all transformations were exact, the reconstructed vector satisfies
\[
 v_1'=kv_2,\qquad
 v_2'=-kv_1+4\rho^{-1}\cos(2ks+\delta)v_1+g/k.
\]
Thus the scalar functions belong to $H^2(0,\infty)$ and solve their
asserted equations. In particular
\[
 \begin{aligned}
 \norm{G_\delta g}_{H^2(0,\infty)}
 &\le C\bigl(\norm{x_{g,\delta}}_{L^2((0,\infty);\R^2)}
                +\norm{g}_{L^2(0,\infty)}\bigr)\\
 &\le\frac CS\norm{g}_{Y_2((0,\infty))},
 \end{aligned}
\]
where $\norm{g}_{L^2(0,\infty)}
\le S^{-2}\norm{g}_{Y_2((0,\infty))}$ was used.
At zero, $R(0)=H_\delta(0)=I$, proving \eqref{eq:normalization}.

Every operator in these formulas is continuous and periodic in $\delta$,
while $e_{\delta+2\pi}=-e_\delta$. The claimed continuity and
antiperiodicity follow, including in the $H^2(0,\infty)$ norm by the exact
reconstructed equations.
\end{proof}

\section{Phase selection and the weak eigenvalue identity}

Take $S$ larger than all the preceding fixed bounds.
By \eqref{eq:T-bound} and Lemma~\ref{lem:scalar},
uniformly in $\delta$ we have
\[
 \norm{G_\delta\Tcal_\delta}_{H^2(0,\infty)\to H^2(0,\infty)}
 \le C/S<1.
\]
Thus
\begin{equation}\label{eq:final-fixedpoint}
 y_\delta=(I-G_\delta\Tcal_\delta)^{-1}h_\delta
\end{equation}
is well defined in $H^2(0,\infty)$ and satisfies \eqref{eq:full-scalar}.
Its normalization remains $\ell_\delta(y_\delta)=1/S$, so it is nonzero.
It is real, continuous in $\delta$, and $y_{\delta+2\pi}=-y_\delta$.
Set
\[
 u_\delta=\alpha_\delta^{-1/2}y_\delta,\qquad
 z_\delta=z_\delta[u_\delta],\qquad
 \psi_\delta=u_\delta\chi+z_\delta.
\]
These functions are real and antiperiodic in $\delta$;
$u_\delta$ is continuous in $H^2(0,\infty)$, and $z_\delta,\psi_\delta$
are continuous in $H^2(\D)$. The half-strip function $\psi_\delta$ solves
$\Lop_\delta\psi_\delta=0$ in $L^2(\D)$ and has transverse Dirichlet
boundary values. Its closed component satisfies $z_{\delta,s}(0)=0$.

The real function
\begin{equation}\label{eq:phase-functional}
 \mathcal B_S(\delta)=u_\delta'(0)
 =\alpha_\delta(0)^{-1/2}
   \left[y_\delta'(0)
       -\frac{\alpha_\delta'(0)}{2\alpha_\delta(0)}y_\delta(0)\right]
\end{equation}
is continuous, since derivative trace is continuous on $H^2(0,\infty)$.
It satisfies $\mathcal B_S(\delta+2\pi)=-\mathcal B_S(\delta)$.
The intermediate value theorem gives $\delta_*\in[0,2\pi]$ with
$\mathcal B_S(\delta_*)=0$. For definiteness, select the least zero in
this compact interval. At this phase
\begin{equation}\label{eq:neumann}
 \partial_s\psi_{\delta_*}(0,t)=0.
\end{equation}
Also $\psi_{\delta_*}\ne0$: its $P_{\D}$ component is
$\alpha_{\delta_*}^{-1/2}y_{\delta_*}\chi$ and
$\ell_{\delta_*}(y_{\delta_*})=1/S$.

Extend $b_{\delta_*}$ and $\psi_{\delta_*}$ evenly to the full strip.
The coefficient $b$ is continuous and globally Lipschitz; a derivative
jump at zero is permitted by the admissible class.
Equation \eqref{eq:neumann} ensures that the even extension of $\psi$
belongs to $H^2(\R\times J)$ and that its longitudinal flux
$f^{-1}\partial_s\psi$ has zero trace at the joining section.
In particular $\psi\in H_0^1(\R\times J)$.

For any smooth compactly supported test function with zero transverse
boundary values, integration by parts on both half-strips gives
\begin{equation}\label{eq:weak}
 \int_{\R\times J}
 \left(f^{-1}\psi_s\overline{\varphi_s}
                  +f\psi_t\overline{\varphi_t}\right)\,ds\,dt
 =\lambda_*\int_{\R\times J}f\psi\overline\varphi\,ds\,dt.
\end{equation}
The joining boundary terms vanish by \eqref{eq:neumann}.
Since $f$ and $f^{-1}$ are bounded, both sides of \eqref{eq:weak}
are continuous in the $H^1(\R\times J)$ norm of $\varphi$.
Density proves this identity for every $\varphi\in H_0^1(\R\times J)$.
The representation theorem for the defining form gives
$\psi\in D(H_q)$ and $H_q\psi=\lambda_*\psi$.

Finally, the profile \eqref{eq:profile} is everywhere positive and
\[
 0<q_{S,\delta_*}(s)^2\le\frac{24k}{S+|s|}.
\]
Increasing $S_0$ to at least $12k$ proves $\norm q_\infty\le\sqrt2$
and uniform decay at both ends. Thus $q\in\mathscr Q_1$.
The essential-spectrum theorem of \cite[Theorem~3.1]{KZ} gives
$\sigma_{\rm ess}(H_q)=[E_1,\infty)$, so the eigenvalue is embedded.
The cited source also supplies this essential-spectrum statement for
the Lipschitz profile constructed here. For completeness, $f_{q_a}(s,t)=f_q(s/a,t/a)$, and
\[
 (U_a\psi)(s,t)=a^{-1}\psi(s/a,t/a)
\]
is a unitary map between the two weighted $L^2$ spaces.
Changing variables in the forms gives
$h_{q_a}[U_a\psi]=a^{-2}h_q[\psi]$, hence
$H_{q_a}U_a=a^{-2}U_aH_q$. This proves the assertion
for every $a>0$ and completes the proof of Theorem~\ref{thm:main}.
\qed

\clearpage
\section*{Relation to the source and verification scope}

For any ambient dimension $n+1\ge3$, choose
\[
 \Gamma(s)=(s,0,\ldots,0),\qquad
 \theta(s)=\int_0^s q(r)\,dr,\qquad
 \Theta(s)=(\cos\theta(s),\sin\theta(s),0,\ldots,0)
\]
in the constant relatively parallel normal frame. The curve is unit speed
and smooth, $\Theta$ is a Lipschitz unit vector,
$|\Theta'|=q$ almost everywhere, and
$k_{\rm curve}\cdot\Theta=0$. The immersion is
\[
 \Phi(s,t)=(s,t\cos\theta(s),t\sin\theta(s),0,\ldots,0).
\]
Its metric is exactly $\operatorname{diag}(1+t^2q(s)^2,1)$.
It is injective: its first coordinate determines $s$, and at fixed $s$
its unit transverse direction determines $t$. Thus it also gives a
genuine embedded ribbon, although the repository permits an abstract
immersed strip. Assumption~1 in \cite{KZ} holds.
Uniform decay gives its condition (3.1), and the amplitude bound is
its condition (5.1). Theorem~3.1 and Corollary~5.3 give the essential
spectrum; Theorem~5.1 supplies the stated Hardy inequality.
The construction answers the point-spectrum part of Open Problem~5.6
negatively for the full class in AIM problem 1007.

The dated submission and literature audit accompanying this paper found
no matching earlier resolution in the records checked on 3 November 2026.
That bounded search does not establish publication priority. The existence argument above uses only elementary
Fourier/Sobolev theory, convergent Neumann series, the intermediate
value theorem, the form representation theorem, and the
essential-spectrum theorem in \cite{KZ}.

\begin{thebibliography}{1}
\bibitem{KZ}
D. Krej\v ci\v r\'ik and K. Zahradov\'a,
\emph{Quantum strips in higher dimensions},
Operators and Matrices \textbf{14} (2020), 635--665.
\href{https://doi.org/10.7153/oam-2020-14-41}{doi:10.7153/oam-2020-14-41}.
\href{https://files.ele-math.com/articles/oam-14-41.pdf}{Published paper};
\href{https://arxiv.org/abs/1906.02431}{arXiv:1906.02431}.
\end{thebibliography}
\end{document}
